\documentclass{article}
\usepackage[margin=1cm]{geometry}
\usepackage{natbib}

\title{Corpus Callosum Dysgenesis and Autistic Traits: The Role of Brain Connectivity}
\date{}


\begin{document}

\maketitle

\section{Introduction}

The corpus callosum is the major white-matter structure connecting the two cerebral hemispheres and is commonly regarded as the principal route for sharing sensory information and coordinating processing across the hemispheres \citep{Banich1998}. Corpus callosum dysgenesis (CCD), occurring in approximately 1 in 4,000 live births, refers to a spectrum of developmental anomalies ranging from partial underdevelopment to complete agenesis \citep{Paul2007}. No single gene or combination of genes has been consistently identified across patients, and the condition does not appear to follow a strongly inherited pattern \citep{Lau2013}. Interestingly, despite the absence or malformation of this major commissural pathway, most individuals with CCD do not exhibit the classic split-brain syndrome, in which the hemispheres function in relative isolation \citep{Brown2019}. This relative preservation of coordinated behavior suggests that alternative pathways and compensatory mechanisms may support interhemispheric integration. At the same time, CCD is associated with highly heterogeneous clinical presentations, ranging from asymptomatic cases to profiles involving psychomotor and communication difficulties, epilepsy, and frequent co-occurring DSM-5 neurodevelopmental conditions, including traits reminiscent of autism spectrum disorder (ASD) \citep{Hantzsch2023}. Given that CCD entails a fundamental alteration of large-scale interhemispheric connectivity, examining the presence and nature of autistic-like traits in this population provides a unique opportunity to clarify how variations in structural connectivity relate to social-communicative phenotypes and to better understand how the developing brain reorganizes when a major commissural structure is absent.

\section{Corpus callosum dysgenesis: development, phenotype, and connectivity}

\subsection{The role of the CC}

The corpus callosum (CC) plays a central role in interhemispheric coordination by regulating how information and processing resources are distributed between the two cerebral hemispheres \citep{Banich1998}. At a basic level, it enables the transfer of sensory and perceptual information across hemispheres, ensuring that inputs from one side of the body or visual field are accessible to the contralateral cortex \citep{Banich1998,Gazzaniga2000}. However, its function extends far beyond simple information transfer. The CC contributes to the organization of cognitive processing by dynamically balancing hemispheric specialization and interhemispheric integration. Regarding specialization, functional lateralization improves efficiency by allowing one hemisphere to specialize in specific operations, reducing redundancy and limiting interhemispheric conduction delays. The CC supports this specialization by regulating interhemispheric interference: callosal projections terminate in cortical layers where they can engage both excitatory pyramidal neurons and inhibitory GABAergic interneurons, allowing activity in one hemisphere to suppress competing activity in the other when necessary. This mechanism enables stable lateralization while preserving flexibility.

At the same time, the CC enables the coordination of processing across hemispheres when task demands exceed the capacity of a single hemisphere, which requires interhemispheric integration. It has been proposed that the CC acts as a dynamic attentional module, allowing one hemisphere to recruit processing resources from the other, thereby supporting parallel computation and increasing overall processing capacity \citep{Banich1998}. Neuropsychological and functional studies support this hypothesis by showing that although each hemisphere can perform simple perceptual or motor operations independently, complex cognitive functions—such as decision-making, contextual integration, or multi-domain processing—require efficient interhemispheric communication and precise temporal coordination \citep{Gazzaniga2000,Hinkley2016,Yuan2020}. Thus, rather than enforcing constant coupling between hemispheres, the CC acts as a regulatory structure that flexibly modulates interhemispheric connectivity, ensuring an optimal trade-off between specialization, parallel processing, and integration.

\subsection{The emergence of DCC}

The corpus callosum develops through a tightly regulated sequence during early brain ontogeny. Cortical projection neurons extend axons toward the midline, where specialized glial structures—such as the glial wedge, midline zipper glia, and indusium griseum—provide guidance cues that enable axons to cross to the contralateral hemisphere and innervate homotopic targets \citep{Tovarmoll2007,Paul2012}. Disruptions in these guidance mechanisms prevent callosal axons from crossing, resulting in partial or complete dysgenesis of the corpus callosum. When axons fail to cross, they usually form longitudinal Probst bundles, which remain ipsilateral and are consistently observed in humans with DCC \citep{Tovarmoll2007,Lazarev2016}. Some individuals also develop sigmoid bundles, connecting non-homologous regions \citep{Tovarmoll2007}. Despite this profound disruption of the principal interhemispheric pathway, the absence of the classical split-brain syndrome observed after surgical callosotomy in adulthood suggests that functional integration across hemispheres is preserved through alternative developmental trajectories. Rather than producing a simple loss of interhemispheric communication, early callosal maldevelopment appears to trigger large-scale reorganization of neural circuitry, leading to a distinct but viable configuration of brain connectivity.

\subsection{Consequences of corpus callosum absence on brain circuitry}

Resting-state fMRI studies have consistently shown that the large-scale organization of canonical resting-state networks is broadly preserved in individuals with dysgenesis of the corpus callosum (DCC), despite the absence or malformation of the principal interhemispheric commissural pathway. Using independent component analysis, \citep{Tyszka2011} demonstrated that bilaterally organized resting-state networks—including the default mode, visual, and sensorimotor systems—can be detected even in cases of complete agenesis, suggesting that the emergence of recognizable functional networks does not strictly require direct callosal fibers \citep{Tyszka2011}. \citep{Owen2013} further reported a modular organization comparable to that observed in controls, indicating that large-scale community structure can develop under markedly atypical structural constraints \citep{Owen2013}. These findings challenge simplistic disconnection models and suggest that large-scale functional architecture can emerge through alternative developmental mechanisms. However, more fine-grained quantitative analyses indicate that this preservation is only partial. Homotopic interhemispheric functional connectivity is consistently reduced across cortical regions, particularly within default mode and fronto-parietal systems \citep{Hinkley2012}. At the same time, several reports describe a relative increase in intrahemispheric connectivity and stronger local clustering within each hemisphere, suggesting a shift toward more segregated, hemisphere-specific processing \citep{Yuan2020,Shi2021}. In this configuration, functional communication appears to rely more heavily on within-hemisphere integration, with diminished cross-hemispheric synchronization. Importantly, this reorganization is heterogeneous and appears sensitive to the degree of callosal malformation, suggesting that while canonical resting-state networks are detectable in DCC, their dynamic integration and interhemispheric synchronization are substantially modified rather than fully preserved.

Such a shift toward within-hemisphere processing has led to the hypothesis that information may become bilaterally represented within each hemisphere, thereby reducing the need for direct interhemispheric transfer. If this were the case, functions that typically require callosal communication should remain intact because each hemisphere would contain redundant representations. For example, under typical organization, tactile input from the left hand is processed in the right somatosensory cortex and must be transferred via the corpus callosum to left-hemisphere language areas for verbal identification; true bilateral representation would therefore predict preserved left-hand naming even in the absence of callosal transfer. \citep{Monteiro2019} directly tested this prediction in the cheirosensory domain. Although individuals with callosal dysgenesis showed intact left-hand naming, neural responses remained strongly lateralized, with no evidence of ipsilateral somatosensory activation \citep{Monteiro2019}. Thus, preserved performance did not reflect genuine bilateral duplication of function. Instead, the persistence of hemispheric specialization suggests compensatory mechanisms, such as indirect interhemispheric pathways or alternative functional reorganization that bypass canonical callosal transfer.

From a developmental perspective, failed callosal crossing does not eliminate axonal growth but redirects it, giving rise to alternative structural configurations. Callosal axons that do not reach the contralateral hemisphere instead form ipsilateral longitudinal bundles known as Probst bundles, which preserve the mediolateral and rostrocaudal topographic gradients characteristic of typical callosal projections \citep{Tovarmoll2007,Wahl2009}. This preserved topography indicates that these fibers retain aspects of their original cortical targeting rules despite failing to cross the midline, suggesting an organized developmental reallocation rather than random misrouting. Diffusion-based measures further indicate that individuals with more prominent Probst bundles tend to exhibit stronger intrahemispheric structural connectivity, paralleling functional findings of increased within-hemisphere connectivity and reduced interhemispheric integration \citep{Paul2012,Yuan2020,Shi2021}. Together, these observations support the view that Probst bundles reinforce longitudinal intrahemispheric communication within an overall network architecture biased toward segregation. In contrast, sigmoid bundles—aberrant interhemispheric tracts connecting non-homologous cortical regions—are preferentially observed in the most severe forms of callosal malformation and are generally interpreted as maladaptive miswiring rather than adaptive reorganization \citep{Owen2013}.

Beyond intrahemispheric reorganization, alternative commissural pathways may contribute to residual interhemispheric transfer in DCC. The anterior commissure, which normally connects anterior temporal regions, is frequently enlarged in DCC \citep{Hetts2006}, and its microstructural integrity has been associated with better cognitive and attentional performance \citep{Siffredi2019}, suggesting partial functional compensation. By contrast, the posterior commissure in the typical brain primarily subserves midbrain and pretectal circuitry, including components of the pupillary light reflex, rather than large-scale cortico-cortical integration \citep{Buttner2014}. Its involvement in cortical communication is therefore not characteristic of normal organization. Nevertheless, diffusion tractography studies in DCC have identified anomalous interhemispheric bundles bypassing the absent corpus callosum by crossing through dorsal midbrain or ventral forebrain regions. \citet{Tovarmoll2014} described ectopic cortico-cortical tracts connecting bilateral parietal areas via noncanonical trajectories involving the anterior or posterior commissural region \citep{Tovarmoll2014}. These bundles were absent in controls and followed organized, reproducible pathways, indicating structured developmental rerouting rather than random miswiring.


An alternative, but not mutually exclusive, interpretation is that subcortical circuits play a central role in preserving bilateral functional connectivity, particularly within sensory systems. In the healthy brain, thalamocortical projections relay sensory information via largely symmetrical bilateral connections to primary sensory cortices. In DCC, resting-state fMRI shows that while interhemispheric connectivity is reduced for associative hubs, it is relatively preserved for primary sensory networks, consistent with compensation via subcortical routes that help enforce bilaterally coordinated activity \citep{Owen2013}. However, transmission through all these alternative pathways is necessarily slower and less temporally precise than through callosal fibers, as they often involve longer and potentially polysynaptic routes with fewer fast, densely myelinated cortico-cortical fibers. This is consistent with electrophysiological and neuroimaging evidence of delayed interhemispheric transfer and reduced synchronization in DCC, particularly for higher-order associative networks \citep{Paul2012,Hinkley2016}.

Overall, all these findings support the hypothesis that, in the absence of the corpus callosum, the brain recruits a heterogeneous set of structural and subcortical adaptations to preserve interhemispheric connectivity, with the relative contribution of each pathway varying across individuals due to developmental plasticity and anatomical constraints. This reorganization provides an essential framework for interpreting the cognitive and behavioral characteristics associated with dysgenesis of the corpus callosum.

\subsection{Consequences of corpus callosum absence on behavior}

According to \citet{Brown2019}, the core syndrome of dysgenesis of the corpus callosum (DCC) includes three main aspects: (1) reduced interhemispheric transfer of sensory-motor information, (2) slowed cognitive processing, and (3) deficits in complex reasoning and novel problem-solving.

First, individuals with DCC show reduced interhemispheric transfer of sensory-motor and tactile-spatial information. Experimental paradigms requiring cross-hemispheric integration—such as tasks in which tactile input from one hand must be verbally identified—consistently reveal increased interhemispheric transfer time and reduced efficiency compared to neurotypical controls \citep{Brown1999,Paul2007}. This reduced efficiency is also evident in motor coordination tasks: on the Bimanual Coordination Test (BCT), individuals with DCC are slower and less accurate when coordinating bilateral movements \citep{Mueller2009}. Notably, their performance resembles that of neurotypical children whose corpus callosum is not yet fully myelinated \citep{Marion2003}, suggesting that DCC reflects a persistent state of diminished interhemispheric integration. Accordingly, transfer is generally sufficient for simple unimanual tasks but becomes less efficient when precise bilateral synchronization or integration of lateralized sensory input is required, consistent with the role of the corpus callosum in coordinating activity across hemispheres.

Second, individuals with DCC exhibit slowed cognitive processing, which has been identified as a primary contributor to executive control difficulties \citep{Brown2019,Marco2012}. This slowing is typically reflected in increased reaction times across a range of cognitive tasks, even when accuracy is preserved. It appears to affect processing speed broadly rather than a single domain, and may limit the efficient coordination of multiple cognitive operations, thereby impacting executive functions such as inhibition, cognitive flexibility, and working memory.

Third, individuals with DCC show impairments in complex reasoning and novel problem-solving. Tasks requiring abstraction, flexible strategy formation, generalization, or integration of multiple conceptual elements are particularly affected, whereas overlearned or routine tasks remain relatively preserved \citep{Brown2019,Erickson2013,Cherbuin2005}. This profile highlights the importance of interhemispheric communication for high-level cognitive integration beyond basic sensory-motor coordination.

As a consequence of these primary deficits, individuals with DCC often show difficulties in social and emotional functioning. They may have trouble recognizing complex emotions, understanding social cues, and integrating cognitive and emotional information in social contexts \citep{Paul2021,Bridgman2014,Anderson2017}. For example, while basic emotional experience is typically intact, the ability to interpret nuanced emotional expressions or predict others’ behavior is impaired, reflecting deficits in higher-order social cognition rather than in perception \citep{Paul2021}. Consistent with this distinction, \citet{Symington2010} reported impairments in theory of mind and in the interpretation of complex social intentions, whereas more basic aspects of social perception were relatively preserved \citep{Symington2010}. Difficulties also extend to adaptive social behavior: individuals may apply social rules rigidly, struggle with flexible decision-making, or fail to use context to guide responses \citep{Brown2021,Mangum2021,Marco2012}. These impairments are particularly pronounced in complex or novel situations, where effective interhemispheric communication would normally support the integration of multiple informational streams, and are often accompanied by reduced metacognitive awareness of one’s own difficulties \citep{Mangum2018,Mangum2021}.

These social and cognitive characteristics overlap with traits commonly described in autism spectrum conditions, including difficulties in social reciprocity, rigid behavioral patterns, and reduced use of contextual information. Although DCC and autism are distinct conditions, this convergence raises the possibility that alterations in large-scale neural connectivity—central in DCC and widely discussed in autism research—may contribute to partially similar behavioral profiles. To examine this potential convergence more closely, we next consider the behavioral phenotype of autism and the neural mechanisms proposed to underlie it.

\section{Description of autism}

\subsection{Behavioral description}

Autism, or Autism Spectrum Disorder (ASD), is a neurodevelopmental condition affecting about 3\% of the general population \citep{Shaw2025}. Autistic phenotypes are highly heterogeneous, ranging from severe presentations to more subtle forms, sometimes described as the broad autism phenotype \citep{Wheelwright2010}. There is no biological test or neuroimaging marker to confirm the diagnosis; it is determined entirely through behavioral observation. Twin studies indicate a strong genetic contribution \citep{Tick2016}, although no single gene is sufficient to account for the condition, suggesting complex and multifactorial mechanisms. According to the Diagnostic and Statistical Manual of Mental Disorders, Fifth Edition \citep{Harm2013}, the diagnosis is based on two core criteria: (A) persistent deficits in social communication and social interaction across contexts, and (B) restricted and repetitive patterns of behavior, interests, or activities. Criterion A includes difficulties in social-emotional reciprocity, impaired understanding and use of social cues (such as eye contact, facial expressions, gestures, and prosody), and challenges in developing, maintaining, and understanding relationships. Criterion B encompasses repetitive motor movements, insistence on sameness and inflexible adherence to routines, highly restricted or fixated interests, and hyper- or hypo-reactivity to sensory input. Symptoms must be present early in development and cause clinically significant impairment. In addition to the core diagnostic domains, a range of associated features are consistently reported in ASD, including motor coordination difficulties \citep{Fournier2010}, executive dysfunction \citep{Hill2004}, and elevated anxiety \citep{Van2011}. Additionally, attentional processes appear atypical in ASD, in particular with respect to impaired attentional disengagement \citep{Landry2004}, reduced cognitive flexibility and shifting \citep{Hill2004}, and altered allocation of attention to socially salient stimuli \citep{Sasson2014}. Although not included in the formal diagnostic criteria, these dimensions contribute substantially to the autistic phenotype.

Several classical cognitive theories of autism, primarily developed in the 1980s and 1990s, have attempted to account for the diverse manifestations of autism within a single explanatory framework. Executive function theories attribute autistic symptoms to core difficulties in cognitive control, such as planning, flexibility, and inhibition \citep{Hill2004}, while social cognition theories focus on difficulties in understanding others’ emotions and mental states, reflected in altered processing of socially relevant stimuli such as faces or voices \citep{Frith1989}. The Weak central coherence account emphasizes a general bias toward local over global processing, both at the perceptual and conceptual level (focusing on details rather than a meaningful whole), leading in particular to reduced integration of contextual information \citep{Frith1989}. However, accumulating evidence indicates that the putative core deficit proposed by each account is neither systematic nor specific. For instance, a detail-focused or local processing bias is not consistently observed across autistic individuals and can coexist with intact global processing \citep{Mottron2006}. Similarly, executive function impairments, often invoked to explain rigidity and repetitive behavior, are also prominent in other neurodevelopmental conditions such as ADHD, in which social deficits are not a defining feature \citep{Booth2011}. These observations challenge single-deficit models and suggest that no isolated cognitive mechanism fully captures the breadth and heterogeneity of the autistic phenotype.

In response to these limitations, more recent accounts have proposed that autistic characteristics arise from the interaction of multiple non-specific cognitive constraints, which become apparent primarily under conditions of high informational and integrative demands \citep{Benyosef2017}. In this view, difficulties emerge most clearly when tasks require the integration of multiple cues into higher-order representations, rather than during isolated executive, social, or perceptual operations. Empirical findings support this perspective: studies report largely preserved accuracy for basic social, perceptual, and contextual processing in ASD, accompanied by slower responses, but selective impairments when task complexity is maximal—particularly when successful performance depends on the integration of global information or complex social signals \citep{Benyosef2017}. Importantly, this complexity-based account does not posit a single primary cognitive deficit; rather, it characterizes autism as a condition in which integrative processing becomes constrained under increasing computational demands.

Recent alternative, not necessarily incompatible accounts have increasingly proposed that sensory hypersensitivity, together with consequential alterations in attentional regulation, may constitute a primary organizing dimension of autism, rather than secondary features \citep{Kinsbourne1991,Liss2006}. Within this perspective, chronic sensory overreactivity is thought to bias the system toward overselective attention as a compensatory mechanism aimed at managing heightened arousal and reducing overwhelming input. Overselective attention would be characterized by intense focus on a restricted subset of stimuli and difficulty disengaging or shifting \citep{Liss2006}. Such narrowing of attentional scope can in turn promote perseverative and repetitive behaviors, which may function as regulatory strategies that reduce environmental unpredictability and stabilize arousal levels \citep{Kinsbourne1991,Liss2006}. Because social interaction requires rapid, flexible coordination of multiple dynamic cues (e.g., gaze, facial expression, prosody, contextual information), an overfocused and inflexible attentional style may also disproportionately affect social communication. In this framework, sensory hypersensitivity, attentional rigidity, repetitive behavior, and social-communicative difficulties are conceptualized as interrelated consequences of a common alteration in arousal and attentional allocation, suggesting that this sensory–attentional dimension may represent a more proximal mechanism than domain-specific social or executive deficits.

\subsection{Neuro correlates of autism}

Both the complexity account and the sensory–attentional account are coherent with converging neurobiological findings. Structural and functional neuroimaging studies report a coordinated pattern of heightened local connectivity alongside reduced long-distance integration, rather than a simple global disconnection \citep{Just2007,Belmonte2004}. Sensory and perceptual cortices appear to be locally overconnected, producing broadly amplified and less selective neural responses that more readily exceed activation thresholds. Such reduced selectivity blurs the distinction between target and competing inputs, effectively lowering signal-to-noise discrimination at early processing stages. As a result, downstream associative regions may receive degraded or poorly structured input and show reduced or inefficient activation—not necessarily due to intrinsic dysfunction, but because of atypical early-stage filtering \citep{Belmonte2004}. These connectivity patterns are consistent with models of altered excitation–inhibition balance and increased neural noise \citep{Hussman2001}, as well as post-mortem observations of atypical cortical minicolumn organization suggestive of reduced inhibitory surround and enhanced local coupling \citep{Casanova2002}.

This connectivity profile is consistently observed across imaging modalities. Structural MRI studies report region-specific cortical enlargement, particularly in parietal and sensory regions, suggesting disproportionate local cortical development \citep{Piven1997,Courchesne2001}. Diffusion MRI studies reveal an overrepresentation of short-range fibers alongside reduced integrity of long-range white-matter tracts, including callosal pathways, reflected in lower fractional anisotropy and increased radial diffusivity in a subset of autistic individuals \citep{Alexander2007,Shukla2010,Casanova2002}. Functional MRI findings further support this pattern, showing preserved or enhanced local synchronization within regions but reduced functional coupling between distant nodes of large-scale networks involved in executive, social, and integrative processing \citep{Just2007,Mason2008}. Together, these converging results point to a connectivity profile dominated by strong local coupling and weakened large-scale integration, with direct implications for autistic behavioral phenotypes.

A plausible neurodevelopmental origin of altered connectivity in autism involves early disruptions of inhibitory signaling mediated by γ-aminobutyric acid (GABA), the brain’s main inhibitory neurotransmitter \citep{Rubenstein2003,Belmonte2004}. During brain development, GABAergic neurons limit how strongly and how broadly neurons activate together, thereby shaping activity-dependent synaptic selection and elimination \citep{Hensch2004,BenAri2002}. In the cerebellum, this inhibitory role is primarily carried by Purkinje cells, which are GABAergic neurons and provide the sole output of the cerebellar cortex. Post-mortem and imaging studies report reduced cerebellar volume and a selective loss of Purkinje cells in autism \citep{Courchesne1994,Fatemi2012}. Reduced inhibitory regulation during early development allows broader and less specific patterns of neural co-activation, increasing the likelihood that weak or poorly specified synaptic connections are retained. This excessive stabilization of synaptic connections is expected to increase neural tissue volume, providing a mechanistic account of early brain overgrowth observed in autism \citep{Courchesne2001,Hazlett2005}. As development continues, synaptic elimination operates on this inefficiently organized network, yielding atypical pruning and persistent alterations in large-scale functional connectivity.

Together, these findings suggest that altered excitation–inhibition balance—evidenced by GABAergic disruption, Purkinje cell abnormalities, and a pattern of increased short-range alongside reduced long-range connectivity—offers a unified framework linking sensory, attentional, and integrative accounts of autism. Locally amplified and less selective processing provides a plausible basis for sensory hypersensitivity and heightened reactivity, while reduced long-distance coordination constrains the integration of distributed information, particularly under conditions requiring complex, multi-cue processing. Within this architecture, behavioral rigidity and restricted interests can be understood as stabilizing responses to intensified and poorly filtered input, and social-communicative difficulties as emerging when large-scale coordination is required to flexibly combine perceptual, contextual, and affective signals.

\subsubsection{CC in autism}

Considering the overlap between ASD and DCC symptoms, it appears theoretically informative to examine whether the corpus callosum in autism shows structural or microstructural particularities that could contribute to altered large-scale integration. In fact, alterations of the corpus callosum (CC) provide a structural substrate for the broader pattern of enhanced local processing combined with reduced long-range integration. Early structural MRI studies reported reduced callosal size—particularly when controlling for overall brain volume—alongside relatively enlarged cortical regions, especially in parietal areas \citep{Piven1997}. Subsequent diffusion imaging studies have consistently identified reductions in callosal microstructural integrity, most prominently in the genu and splenium, reflected in decreased fractional anisotropy, reduced fiber density, or altered fixel-based metrics suggestive of diminished axonal density \citep{Just2007,Booth2011,Frazier2009,Dimond2019}. These alterations disproportionately affect anterior and posterior segments supporting fronto-parietal and temporo-parietal interactions—networks critically involved in executive control, theory of mind, and multimodal integration.

Functional imaging further indicates that, although activation within individual cortical regions is often preserved, connectivity between distributed regions—particularly across hemispheres—is reduced \citep{Just2007}. Reduced callosal integrity is associated with decreased interhemispheric functional connectivity and correlates with behavioral measures of social impairment; for example, lower splenial fiber density has been shown to predict greater social symptom severity \citep{Dimond2019}. Anatomical evidence of a higher proportion of short-range relative to long-range fibers \citep{Casanova2002} suggests that CC alterations form the interhemispheric expression of a broader bias toward local over distributed connectivity. Importantly, these callosal differences are not uniform: group-level reductions coexist with substantial overlap between autistic and control individuals, and milder but more generalized white matter alterations beyond the CC \citep{Dimond2019}, indicating that interhemispheric under-connectivity represents a graded and heterogeneous mechanism contributing to symptom variability rather than a categorical structural deficit \citep{Paul2007,Frazier2009}.

\section{Comparison between ASD and DCC}

Given the substantial overlap in clinical features between autism and DCC, it is likely that some aspects of altered connectivity are shared. Although autism does not involve a complete absence of the corpus callosum, both conditions are marked by disruptions in large-scale integration: DCC through the absence or malformation of corpus callosum, and autism through reduced long-range connectivity. In both cases, tasks requiring integration of distributed information—particularly across hemispheres or distant cortical regions—appear especially vulnerable. Moreover, while autism has been described in terms of increased short-range overconnectivity, DCC is likewise associated with atypical strengthening of alternative intrahemispheric or subcortical pathways that may partially compensate for callosal absence, potentially resulting in a qualitatively comparable imbalance between local and global connectivity. Despite distinct developmental origins, partially convergent patterns of network organization may therefore contribute to similar cognitive and behavioral outcomes, highlighting the need for direct comparative investigation.




\section{Behavioral and connectivity comparison between ASD and DCC}

\subsection{Autistic-like symptoms in DCC}

Most studies directly comparing DCC with ASD focus on behavioral aspects. Across comparative studies, individuals with DCC show elevations in social interaction difficulties, theory of mind, emotion understanding, and pragmatic communication relative to neurotypical controls, though typically to a lesser extent than individuals with ASD \citep{Badaruddin2007,Lau2013}. Executive and attentional differences follow a similar pattern: impairments in cognitive flexibility, sustained attention, and susceptibility to interference are consistently observed in ASD and are also present in DCC, but generally at an intermediate level between controls and autistic participants \citep{Badaruddin2007,Paul2014}. Notably, measures tapping socially embedded attentional demands—such as rapid integration of contextual or prosodic cues—tend to accentuate this gradient pattern. Across domains, DCC participants often occupy a position between neurotypical and ASD groups, suggesting partial but not complete phenotypic overlap. Moreover, social difficulties in DCC may emerge later in development or become more apparent in complex real-world contexts, which may contribute to under-recognition of overlapping autistic features \citep{Booth2018}.

The clearest divergence between DCC and ASD emerges in studies that directly compare the restricted and repetitive behaviors and interests (RRBI) domain. Whereas ASD is characterized by circumscribed interests, stereotyped behaviors, insistence on sameness, and atypical sensory responsivity, these features are generally less pronounced in DCC. In comparative samples, repetitive–restrictive behaviors tend to be the least affected domain in DCC relative to ASD \citep{Badaruddin2007}, with lower frequencies of stereotyped movements and intense restricted interests reported on standardized measures \citep{Lau2013}. Although some individuals with DCC may exhibit preference for routine or mild perseverative tendencies—often discussed in relation to executive functioning—these behaviors are typically less pervasive and less central to the behavioral profile than in ASD. Evidence for a detail-focused perceptual style is likewise limited in DCC cohorts \citep{Booth2018}. With respect to sensory features, most early comparative studies did not report robust elevations in hypersensitivity in DCC; however, more recent work suggests that atypical sensory responsivity may be present in a subset of individuals with callosal anomalies, albeit with a different profile and lower severity than in ASD \citep{Demopoulos2015}. Overall, direct behavioral comparisons indicate a relative attenuation of RRBI and sensory-driven features in DCC compared to ASD, even when social-communicative differences are clearly observed.

Taken together, DCC mirrors a substantial portion of the social-cognitive and executive components of ASD but does not recapitulate the full phenotype. Its profile is both quantitatively milder and qualitatively incomplete, supporting the view that the autistic phenotype reflects separable dimensions—only some of which appear to depend strongly on interhemispheric integration. A parsimonious hypothesis is that the overlap between ASD and DCC is primarily driven by vulnerabilities in long-range connectivity and large-scale integration. In DCC, the absence or malformation of the corpus callosum directly constrains interhemispheric communication; in ASD, converging neuroimaging studies report reduced long-distance coordination and inefficient cross-network integration \citep{Just2007,Belmonte2004}. These shared integrative constraints may account for convergent difficulties in complex social cognition, contextual processing, and executive control. By contrast, features more strongly associated with local circuitry alterations—such as sensory hypersensitivity, heightened reactivity, and pronounced repetitive–restrictive behaviors—appear attenuated or inconsistently expressed in DCC and have been linked in ASD to increased short-range connectivity and excitation–inhibition imbalance \citep{Casanova2002,Hussman2001}. This pattern suggests a dissociation whereby large-scale integrative disruption underlies shared social-cognitive features, whereas additional local circuit alterations contribute to the full expression of RRBI and sensory hyperreactivity in ASD.

An additional important observation concerns the dissociation between structural extent and behavioral severity. Despite representing a more overt anatomical alteration than the callosal reductions typically reported in ASD \citep{Frazier2009,Paul2007}, complete agenesis of the corpus callosum is often associated with a milder behavioral phenotype. Moreover, within DCC, complete agenesis does not systematically result in more severe social-communicative impairments than partial dysgenesis, a pattern attributed to early developmental reorganization and compensatory recruitment of alternative pathways \citep{Tovarmoll2014,Brown2019}. This discrepancy suggests that behavioral outcomes cannot be inferred from structural disruption alone but depend on how large-scale networks reorganize over development.

\subsection{Behavioral correlations}

Most direct comparative studies have focused on differences in symptom severity rather than on relationships among symptoms. However, examining correlations between behavioral dimensions provides a more informative approach, as stable associations are less likely to arise incidentally and may reveal underlying organizational principles.

In autism, social-communication impairments are closely intertwined with attentional regulation, processing speed, and executive flexibility. Greater social symptom severity correlates with poorer joint attention, reduced adaptive communication, and diminished pragmatic language skills \citep{Mundy2007,TagerFlusberg2006}. Deficits in attentional orienting toward socially relevant stimuli are associated with increased social impairment \citep{Keehn2013}, and greater intra-individual reaction time variability correlates with communication difficulties and reduced adaptive functioning \citep{Dinstein2012}. Cognitive flexibility deficits correlate with impairments in social reciprocity and pragmatic language \citep{Lopez2005}, while slower processing speed is associated with weaker executive control and poorer conversational reciprocity \citep{Marco2012}. Together, these correlations delineate a cluster centered on integrative efficiency and cognitive flexibility, paralleling the behavioral profile described in DCC.

By contrast, autism also exhibits a distinct behavioral subsystem linking sensory hypersensitivity, anxiety, intolerance of uncertainty, and repetitive behaviors. Sensory over-responsivity correlates strongly with repetitive behaviors and predicts anxiety severity \citep{Schulz2019,Green2015}. Intolerance of uncertainty partially mediates this relationship \citep{Wigham2015}, suggesting a cascade from sensory reactivity to anxiety and behavioral rigidity. These associations remain significant even when controlling for social symptoms, indicating a partially dissociable axis. Given that DCC does not consistently reproduce this coupling, this second cluster may reflect a distinct mechanistic architecture specific to ASD.

\subsection{Connectivity--symptom relationships}

To date, no study has directly compared ASD and DCC in terms of connectivity–symptom coupling. However, autism research shows that distinct symptom clusters map onto partially dissociable large-scale network alterations.

In ASD, social-communication impairments are associated with reduced long-range connectivity within and between associative networks. Reduced functional connectivity within the default mode network (DMN) correlates with social impairment severity and poorer theory-of-mind performance \citep{Assaf2010,Elton2016}. Structural integrity of callosal fibers linking bilateral associative regions further predicts social symptom levels \citep{Dimond2019}. Reduced connectivity between cerebellar associative regions and cortical social networks also correlates with social deficits \citep{Ramos2019,Khan2015}, consistent with models of disrupted cerebellar modulation \citep{Belmonte2004}.

Attention and executive processes are embedded within this architecture. Reduced connectivity in the dorsal attention network correlates with attentional variability \citep{Just2007}, while altered cerebellar coupling relates to shifting deficits \citep{Rohr2019}. Executive control network connectivity is associated with modulation of communication, and correlates with callosal structure in ASD \citep{Booth2011}.

In contrast, the sensory–rigidity dimension is associated with a distinct connectivity profile involving salience, thalamic, and attentional systems. Sensory hypersensitivity correlates with increased salience network connectivity and stronger coupling with primary sensory cortices \citep{Green2015,Uddin2013}. Thalamocortical hyperconnectivity correlates with sensory symptoms \citep{Nair2013,Baran2023}, while local neural synchrony is associated with excitation–inhibition imbalance \citep{Orekhova2007}. Repetitive behaviors correlate with salience network hyperconnectivity, reduced executive control connectivity, and altered fronto-striatal circuits \citep{Langen2012}.

Together, these findings delineate two partially dissociable architectures in ASD: a long-range integrative network underlying social-communication functions, and a sensory–salience–rigidity network supporting hypersensitivity and repetitive behavior.

Building on this framework, the present study aims to examine whether similar connectivity–symptom relationships can be identified in DCC. Comparing these structured patterns across conditions may provide a more mechanistic understanding of shared and distinct features than symptom-level comparisons alone.

\bibliographystyle{plainnat}
\bibliography{PRONAS_clean}

\end{document}
